\section{Integration}\label{sec:integration}
In this section, we use the Poisson summation formula (see \cite[Proposition~4.2]{Ohtsuki:QT2016}) to change the summation in \eqref{eq:J_summation} into an integration.
Then by using the saddle point method (see also \cite[Proposition~3.2]{Ohtsuki:QT2016}) to prove Theorem~\ref{thm:main}.
\par
Define
\begin{equation*}
  \psi(z)
  :=
  F(z+\varphi(\xi)/\xi)-F(\varphi(\xi)/\xi)
\end{equation*}
in the region
\begin{equation}\label{eq:region_psi}
  \left\{
    z\in\C
    \Bigm|
    -\frac{\varphi(\xi)}{\xi}\le\Re{z}<1-\frac{\varphi(\xi)}{\xi},
    -1-\frac{\varphi(\xi)}{\xi}<\frac{\Im{z}}{\tan\theta}+\Re{z}<1-\frac{\varphi(\xi)}{\xi}
  \right\}.
\end{equation}
Then since $\psi(0)=0$, $\psi'(0)=0$, and $\psi''(0)=-\xi\sqrt{(2\cosh(\xi)-1)^2-4}$ from \eqref{eq:F'_phi} and \eqref{eq:F''_phi}, $\psi(z)$ is of the form
\begin{equation*}
  \psi(z)
  =
  \frac{-\xi}{2}\sqrt{(2\cosh(\xi)-1)^2-4}\times z^2+a_3z^3+a_4z^4+\cdots
\end{equation*}
in \eqref{eq:region_psi}.
\begin{lem}\label{lem:psi_negative}
If $x\ne0$ is real and satisfies $-\varphi(\xi)/\xi\le x<1-\varphi(\xi)/\xi$, then $\psi(x)<0$.
\end{lem}
\begin{proof}
Since $e^{\xi}+e^{-\xi}-e^{\varphi(\xi)}-e^{-\varphi(\xi)}=1$, we have
\begin{equation}\label{eq:psi'}
\begin{split}
  \psi'(x)
  &=
  \log(e^{\xi}+e^{-\xi}-e^{\xi x+\varphi(\xi)}-e^{-\xi x-\varphi(\xi)})
  \\
  &=
  \log(e^{\varphi(\xi)}+e^{-\varphi(\xi)}-e^{\xi x+\varphi(\xi)}-e^{-\xi x-\varphi(\xi)}+1)
  \\
  &=
  \log
  \left(
    e^{-\varphi(\xi)}(1-e^{\xi x})(e^{2\varphi(\xi)}-e^{-\xi x})+1
  \right)
\end{split}
\end{equation}
from \eqref{eq:F'}.
Therefore we see that $\psi'(x)=0$ if $x=0$ or $-2\varphi(\xi)/\xi$, $\psi'(x)>0$ if $-2\varphi(\xi)/\xi<x<0$, and $\psi'(x)<0$ otherwise.
\par
Therefore for $-\varphi(\xi)/\xi<x<0$, $\psi(x)$ is monotonically increasing and for $0<x<1-\varphi(\xi)/\xi$ it is monotonically decreasing.
So $\psi(x)$ takes its unique maximum $0$ at $x=0$, which shows that $\psi(x)<0$ when $x\ne0$.

\end{proof}
Now we use the following proposition.
\begin{prop}[{\cite[Proposition~4.2]{Ohtsuki:QT2016}}]\label{prop:Poisson}
Let $b_{-}$ and $b_{+}$ be real numbers with $b_{-}<0<b_{+}$.
Put
\begin{align*}
  \Lambda
  &:=
  \left\{
    \frac{k}{N}\Bigm| k\in\Z, b_{-}\le\frac{k}{N}\le b_{+}
  \right\},
  \\
  C
  &:=
  \{t\in\R\mid b_{-}\le t\le b_{+}\},
  \\
  D
  &:=
  \{z\in\C\mid\Re{\psi(z)}<0\}.
\end{align*}
Assume that $\psi(z)$ is a holomorphic function of the form
\begin{equation}\label{eq:f}
  \psi(z)
  =
  az^2+a_3z^3+a_4z^4+\cdots
\end{equation}
with $\Re(a)<0$, defined in a neighborhood $P$ of $0\in\C$ that includes the $\delta_0$-neighborhood $N_{\delta_0}$ of $0$ for $\delta_0>0$.
We choose $P$ so that the region $D\cap P$ has two connected components.
We also assume the following:
\begin{enumerate}
\item
$b_{-}$ and $b_{+}$ are in different components of $D\cap P$ and moreover $\Re{\psi(b_{\pm})}<-\varepsilon_0$ for some $\varepsilon_0>0$,
\item
Both $b_{-}$ and $b_{+}$ are in a connected component of
\begin{equation}\label{eq:R+}
  \{x+y\i\mid x\in[b_{-},b_{+}],y\in[0,\delta_0],\Re{\psi(x+y\i)}<2\pi y\}
\end{equation}
\item
Both $b_{-}$ and $b_{+}$ are in a connected component of
\begin{equation}\label{eq:R-}
  \{x-y\i\mid x\in[b_{-},b_{+}],y\in[0,\delta_0],\Re{\psi(x-y\i)}<2\pi y\}
\end{equation}
\end{enumerate}
Then there exists $\varepsilon>0$, depending on $a$ and $\varepsilon_0$, such that
\begin{equation*}
  \frac{1}{N}
  \sum_{z\in\Lambda}
  e^{N\psi(z)}
  =
  \int_{C}e^{N\psi(z)}\,dz
  +
  O(e^{-N\varepsilon}).
\end{equation*}
\end{prop}
We will show that $\psi(z)$ satisfies the assumptions of Proposition~\ref{prop:Poisson}.
\par
Put $b_{-}:=-\varphi(\xi)/(2\xi)$, $b_{+}:=(1-\varphi(\xi)/\xi)/2$, $\delta_0:=\frac{1}{2}(1-\varphi(\xi)/\xi)\sin\theta$.
Define the following regions:
\begin{align*}
  P
  &:=
  \left\{
    z\in\C
    \mid
    -\frac{\varphi(\xi)}{\xi}<\Re{z}<1-\frac{\varphi(\xi)}{\xi},
    -1-\frac{\varphi(\xi)}{\xi}<\frac{\Im{z}}{\tan\theta}+\Re{z}<1-\frac{\varphi(\xi)}{\xi}
  \right\},
  \\
  D_{-}
  &:=
  \{z\in\C\mid\Re{\psi(z)}<0,\Re{z}<0\},
  \\
  D_{+}
  &:=
  \{z\in\C\mid\Re{\psi(z)}<0,\Re{z}>0\}.
\end{align*}
Note that $P$ is just the region \eqref{eq:region_psi} and so $\psi(z)$ is holomorphic in $P$.
See Figure~\ref{fig:P}.
\begin{figure}[h!]
\includegraphics{picture.eps}
\caption{The parallelogram is the boundary of $P$, the circle is the boundary of $N_{\delta_0}$, the red region is $D_{-}$, and the green region is $D_{+}$.
The cyan rectangle indicates the boundary of $R_{+}\cup R_{-}$} (see below).
\label{fig:P}
\end{figure}
\par
Then, we show the following lemma, from which the assumptions of Proposition~\ref{prop:Poisson} hold.
\begin{lem}\label{lem:Poisson_psi}
We assume that $\sin{\theta}<\frac{\varphi(\xi)}{\xi-\varphi(\xi)}$ and that $\tan\theta<\frac{\pi}{2\xi}$.
Then, the function $\psi(z)$ satisfies the following:
\begin{enumerate}
\item[(i).]
$\psi(z)$ is a holomorphic function of the form \eqref{eq:f} defined in the $\delta_0$-neighborhood $N_{\delta_0}$ of $0\in\C$.
\item[(ii).]
Both $D_{+}\cap P$ and $D_{-}\cap P$ are connected, and $b_{\pm}\in D_{\pm}\cap P$.
Moreover, $\Re\psi(b_{\pm})<-\varepsilon_0$ for some $\varepsilon_0>0$.
\item[(iii).]
Condition (2) in Proposition~\ref{prop:Poisson} is satisfied.
\item[(iv).]
Condition (3) in Proposition~\ref{prop:Poisson} is satisfied.
\end{enumerate}
\end{lem}
\begin{proof}
\quad\par
\begin{enumerate}
\item[(i).]
The distances to the upper side, the lower side, the left side, and the right side of $P$ from the origin $O$ is $(1-\varphi(\xi)/\xi)\sin\theta$, $(1+\varphi(\xi)/\xi)\sin\theta$, $\varphi(\xi)/\xi$ and $1-\varphi(\xi)/\xi$, respectively.
Since $(1-\varphi(\xi)/\xi)\sin\theta<\varphi(\xi)/\xi$ from the assumption $\sin\theta<\frac{\varphi}{\xi-\varphi(\xi)}$, $\delta_0$ is less than these distances.
So $N_{\delta_0}$ is contained in the parallelogram $P$.
Moreover since $\psi''(0)=F''(\varphi(\xi)/\xi)=-\xi\sqrt{(2\cosh(\xi)-1)^2-4}<0$, $\psi(z)$ is of the form \eqref{eq:f} with $a:=-\frac{\xi}{2}\sqrt{(2\cosh(\xi)-1)^2-4}$.
\item[(ii).]
We first show that $b_{\pm}\in D_{\pm}\cap P$.
It is clear that $b_{-}<0$, $b_{+}>0$, and $b_{\pm}\in P$.
The inequality $\Re\psi(b_{\pm})<0$ follows from Lemma~\ref{lem:psi_negative}.
So, we can choose $\varepsilon_0>0$ so that $\Re\psi(b_{\pm})<-\varepsilon_0$.
\par
Next, we will show that for each $x\ne0$ with $-\varphi(\xi)/\xi<x<1-\varphi(\xi)/\xi$, the set $\{y\in\R\mid\Re\psi(x+y\i)<0\}\cap P$ is an open interval containing $0$.
Then, we can see that $D_{\pm}\cap P$ is connected.
\par
If we put
\begin{equation*}
  g(x,y)
  :=
  e^{\xi}+e^{-\xi}-e^{\xi(x+y\i)+\varphi(\xi)}-e^{-\xi(x+y\i)-\varphi(\xi)},
\end{equation*}
then we have
\begin{equation}\label{eq:arg}
  \frac{d}{d\,y}\Re\psi(x+y\i)
  =
  -\arg g(x,y)
\end{equation}
from \eqref{eq:psi'}.
Now we have
\begin{align*}
  \Re g(x,y)
  &=
  2\cosh{\xi}
  -
  2\cosh\left(\xi\left(x+\frac{\varphi(\xi)}{\xi}\right)\right)\cos(\xi y),
  \intertext{and}
  \Im g(x,y)
  &=
  -\sinh\left(\xi\left(x+\frac{\varphi(\xi)}{\xi}\right)\right)\sin(\xi y).
  \notag
\end{align*}
\par
We can see that $\Re g(x,y)$ is positive since $0<x+\varphi(\xi)/\xi<1$.
We will show that $\Im g(x,y)<0$ ($\Im g(x,y)>0$) if $y>0$ ($y<0$, respectively).
If $x+y\i\in P$, then $-2\tan\theta<y<\tan\theta$ (see Figure~\ref{fig:P}).
Together with the assumption $\tan\theta<\frac{\pi}{2\xi}$, we conlude that $|y|\le2\tan\theta<\pi/\xi$, and so $|\xi y|<\pi$.
This implies that $\sin(\xi y)>0$ ($\sin(\xi y)<0$, respectively) if $y>0$ ($y<0$, respectively).
Since $x+\varphi(\xi)/\xi>0$, we see that $\Im g(x,y)<0$ ($\Im g(x,y)>0$) if $y>0$ ($y<0$, respectively).
\par
Therefore $\arg{g(x,y)}$ is positive (negative, respectively) when $y<0$ ($y>0$, respectively).
It follows that if $y>0$ ($y<0$, respectively), $\Re\psi(x+y\i)$ is monotonically increasing (decreasing, respectively) with respect to $y$.
Since $\Re\psi(x)<0$ from Lemma~\ref{lem:psi_negative}, we see that $\{y\in\R\mid\Re\psi(x+y\i)<0\}\cap P$ is an open interval containing $0$.
\par
Therefore, we conclude that $D_{\pm}\cap P$ is connected.
\item[(iii).]
We will show that the region $R_{+}:=[b_{-},b_{+}]\times[0,\delta_0]\setminus\{O\}$ is contained in the region $\{x+y\i\in\C\mid\Re\psi(x+y\i)<2\pi y\}$, which means that the region \eqref{eq:R+} equals $R_{+}$.
Then we can find a path connecting $b_{-}$ and $b_{+}$ in $R_{+}$, proving Condition (2).
\par
In fact, we can show if $b_{-}\le x\le b_{+}$ and $0\le y\le\delta_0$ ($(x,y)\ne(0,0)$), then $\Re\psi(x+y\i)-2\pi y<0$ as follows.
\par
When $y>0$, then $\Im g(x,y)<0$ from the argument in (ii), and so $g(x,y)$ is in the fourth quadrant in the complex plane because $\Re g(x,y)>0$.
Therefore we have $0>\arg g(x,y)>-\pi/2$.
This shows that $\frac{d}{d\,y}\bigl(\Re\psi(x+y\i)-2\pi y\bigr)<0$ when $y>0$ from \eqref{eq:arg}, and so we conclude that $\Re\psi(x+y\i)-2\pi y<0$ since $\Re\psi(x)\le0$ from Lemma~\ref{lem:psi_negative}.
\item[(iv).]
We will show that the region $R_{-}:=[b_{-},b_{+}]\times[-\delta_0,0]\setminus\{O\}$ is in the region $\{x+y\i\in\C\mid\Re\psi(x+y\i)<-2\pi y\}$, which means that the region \eqref{eq:R-} equals $R_{-}\cup\{O\}$.
Then we can find a path connecting $b_{-}$ and $b_{+}$ in $R_{-}$, proving Condition (3).
\par
We will show if $b_{-}\le x\le b_{+}$ and $-\delta_0\le y\le0$ ($(x,y)\ne(0,0)$), then $\Re\psi(x+y\i)+2\pi y<0$.
\par
When $y<0$, then $\Im g(x,y)>0$ from the argument in (ii), and so $g(x,y)$ is in the first quadrant because $\Re g(x,y)>0$.
Therefore we have $0<\arg g(x,y)<\pi/2$.
This shows that $\frac{d}{d\,y}\bigl(\Re\psi(x+y\i)+2\pi y\bigr)>0$ when $y<0$, and so we conclude that $\Re\psi(x+y\i)+2\pi y<0$ since $\Re\psi(x)\le0$ from Lemma~\ref{lem:psi_negative}.
\end{enumerate}
\end{proof}
Now we can prove Theorem~\ref{thm:main}.
\begin{proof}[Proof of Theorem~\ref{thm:main}]
From Proposition~\ref{prop:Poisson}, there exists $\varepsilon>0$ such that
\begin{equation*}
  \frac{1}{N}
  \sum_{-\varphi(\xi)/(2\xi)\le k/N\le(1-\varphi(\xi)/\xi)/2}
  e^{N\psi(k/N)}
  \\
  =
  \int_{-\varphi(\xi)/(2\xi)}^{(1-\varphi(\xi)/\xi)/2}
  e^{N\psi(z)}\,dz
  +
  O\left(e^{-N\varepsilon}\right).
\end{equation*}
Since $f_N(z+\varphi(\xi)/\xi+1/(2N))-S(\xi)/\xi$ uniformly converges to $\psi(z)$, we have
\begin{multline*}
  \frac{1}{N}
  \sum_{-\varphi(\xi)/(2\xi)\le k/N\le(1-\varphi(\xi)/\xi)/2}
  e^{N\left(f_N(k/N+\varphi(\xi)/\xi+1/(2N))-S(\xi)/\xi\right)}
  \\
  =
  \int_{-\varphi(\xi)/(2\xi)}^{(1-\varphi(\xi)/\xi)/2}
  e^{N\psi(z)}\,dz
  +O(e^{-N\varepsilon})
\end{multline*}
from \cite[Remark~4.4]{Ohtsuki:QT2016}.
Putting $l/N:=k/N+\varphi(\xi)/\xi$, we also have
\begin{multline*}
  \sum_{-\varphi(\xi)/(2\xi)\le k/N\le(1-\varphi(\xi)/\xi)/2}
  e^{N\left(f_N\left(k/N+\varphi(\xi)/\xi+1/(2N)\right)-S(\xi)/\xi\right)}
  \\
  =
  e^{-N S(\xi)/\xi}
  \sum_{\varphi(\xi)/(2\xi)\le l/N\le(1+\varphi(\xi)/\xi)/2}
  \exp\left(N\times f_N\left(\frac{2l+1}{2N}\right)\right).
\end{multline*}
Thus we obtain
\begin{multline}\label{eq:Poisson_f_N}
  \sum_{\varphi(\xi)/(2\xi)\le k/N\le(1+\varphi(\xi)/\xi)/2}
  \exp\left(N\times f_N\left(\frac{2k+1}{2N}\right)\right)
  \\
  =
  N\times e^{N S(\xi)/\xi}
  \left(
    \int_{-\varphi(\xi)/(2\xi)}^{(1-\varphi(\xi)/\xi)/2}
    e^{N\psi(z)}\,dz
    +
    O\left(e^{-N\varepsilon}\right)
  \right).
\end{multline}
Now by using the saddle point method \cite[Proposition~3.2]{Ohtsuki:QT2016} (see also \cite[Remark~3.3]{Ohtsuki:QT2016}), we have
\begin{equation}\label{eq:saddle}
  \int_{-\varphi(\xi)/(2\xi)}^{(1-\varphi(\xi)/\xi)/2}
  e^{N\psi(z)}\,dz
  =
  \frac{\sqrt{\pi}}{\sqrt{\frac{\xi}{2}\sqrt{(2\cosh(\xi)-1)^2-4}}\sqrt{N}}
  \left(
    1+O(N^{-1})
  \right).
\end{equation}
Putting $c_{-}:=\varphi(\xi)/(2\xi)$ and $c_{+}:=(1+\varphi(\xi)/\xi)/2$ in \eqref{eq:J_summation}, we finally have
\begin{equation*}
\begin{split}
  &J_N\bigl(E;\exp(\xi/N)\bigr)
  \\
  =&
  \frac{N\times e^{N S(\xi)/\xi}}{2\sinh(\xi/2)}
  \left(
    \int_{-\varphi(\xi)/(2\xi)}^{(1-\varphi(\xi)/\xi)/2}
    e^{N\psi(z)}\,dz
    +O\left(e^{-N\min\{\varepsilon,\delta'\}}\right)
  \right)
  \\
  =&
  \frac{\sqrt{\pi}\sqrt{N} e^{N\times S(\xi)/\xi}}
       {2\sinh(\xi/2)\sqrt{\frac{\xi}{2}\sqrt{(2\cosh(\xi)-1)^2-4}}}
  \left(1+O(N^{-1})\right).
\end{split}
\end{equation*}
where we use \eqref{eq:Poisson_f_N} at the first equality and \eqref{eq:saddle} at the second.
Letting $T(\xi)$ denote $\frac{2}{\sqrt{(2\cosh(\xi)-1)^2-4}}$, we have
\begin{equation*}
  J_N\bigl(E;\exp(\xi/N)\bigr)
  =
  \frac{\sqrt{\pi}}{2\sinh(\xi/2)}
  \sqrt{T(\xi)}
  \sqrt{\frac{N}{\xi}}
  \times
  e^{\frac{N}{\xi}S(\xi)}
  \left(1+O(N^{-1})\right).
\end{equation*}
Since $J_{N}(E;e^{-\xi/N})=J_{N}(E;e^{\xi/N})$ we obtain the required formula.
\end{proof}